\section{Secure Boot}

\begin{frame}{Purpose}
  \begin{itemize}
    \item Start from the following premises:
    \begin{itemize}
      \item the system's software was installed securely
      \item the system enters production and might be exposed to threats
      \item some of those threats will know how to install software
    \end{itemize}
    \item The question is: how do I make sure the software has not been altered?
    \item The system must be able to assess autonomously
  \end{itemize}
\end{frame}

\begin{frame}{Concept}
  \begin{itemize}
    %% As we've already seen, systems boot in stages
    \item Chain of verification of software
    \begin{itemize}
      \item the manufacturer's ROM code verifies the first stage bootloader
      \item the last stage bootloader (e.g. U-Boot) verifies the OS kernel
      \item the kernel, optionally, verifies the userland
    \end{itemize}
    %% Of course, more complicated for userland
    \item Compute the hash of the next stage binary
    \item Use the embedded \textbf{public} key to retrieve the signed hash
    \item Compare re-computed hash to extracted hash
    \item Refuse to hand over control flow if the hashes do not match
  \end{itemize}
\end{frame}

\begin{frame}{Initial bootloader verification}
  \begin{columns}
    \column{0.7\textwidth}
      \includegraphics[width=0.9\textwidth]{slides/security-secure-boot-part-1/phase1.pdf}
    \column{0.3\textwidth}
      \begin{itemize}
        \item The public key's hash is stored in a once-writable memory
        \item The BootROM uses it to validate the embeded key
        \item It uses that key to recompute the hash from the signature
      \end{itemize}
  \end{columns}
\end{frame}

\begin{frame}{Initial bootloader verification}
  \begin{itemize}
    \item Storing the hash instead of the full public key means the storage is
          independent of the signature algorithm.
          It only depends on the hash algorithm.
    \item The BootROM must be trusted. It cannot be modified, as it is stored in
          a read-only memory.
    \item The key hash is stored in a memory that can only be written once.
          This is usually done by burning electrical fuses.
    \item In some implementations, one more fuse must be burned once the
          implementation has been verified to be correct to enable secure boot.
  \end{itemize}
\end{frame}

\begin{frame}{Kernel verification}
  \begin{columns}
    \column{0.7\textwidth}
      \includegraphics[width=0.9\textwidth]{slides/security-secure-boot-part-1/phase2.pdf}
    \column{0.3\textwidth}
      \begin{itemize}
        \item End of the "traditional" secure boot chain
        \item Verify the kernel before execution
        \item Must be supported by the bootloader
      \end{itemize}
  \end{columns}
\end{frame}

\begin{frame}{Kernel verification}
  \begin{itemize}
    \item This stage no longer depends on hardware
    \item The root of trust for this stage is the combined:
    \begin{itemize}
      \item last stage bootloader (e.g. U-Boot proper)
      \item last stage signature public key
    \end{itemize}
    \item $Signature^{-1}$
    \item Both should have been verified in the previous stage
    \item This ensures that the kernel that has been started is the one
          that the possessor of the last stage private key expects.
  \end{itemize}
\end{frame}

\begin{frame}{Kernel verification}
  \begin{itemize}
    \item This leaves some unanswered questions:
    \begin{itemize}
      \item What happens once the kernel starts \textit{init}?
      \item What about the DTB?
      \item What about the command line parameters?
    \end{itemize}
    \item We'll answer these in the next part.
  \end{itemize}
\end{frame}

\begin{frame}{Examples}
  \begin{itemize}
    \item In the x86/PC world:
    \begin{itemize}
        \item usually implemented via \href{https://uefi.org/}{UEFI}.
          %%
        \item Pushed by Microsoft,
              \href{https://blogs.windows.com/windows-insider/2021/08/27/update-on-windows-11-minimum-system-requirements-and-the-pc-health-check-app/}
                   {Windows 11 makes it a requirement}
    \end{itemize}
    \item On Android: \href{https://source.android.com/docs/security/features/verifiedboot/avb}{Android Verified Boot} (AVB)\newline
          This is why "\href{https://source.android.com/docs/security/features/verifiedboot/boot-flow}{bootloader unlocks}" are necessary
    \item In \href{https://docs.u-boot.org/en/latest/usage/fit/verified-boot.html}{U-Boot},
    FIT (Flat Image Tree ) images can be signed.
  \end{itemize}
\end{frame}

\begin{frame}{Threat Model}
  \begin{itemize}
    \item Secure boot is designed to protect against unauthorized modification
          of software. This includes:
    \begin{itemize}
      \item Offline modification of the software by reflashing the boot medium.
            This could be how an attacker gain access.
      \item Runtime rewriting of the software (persistence).
            This is defense-in-depth against an attacker who already has access.
    \end{itemize}
    \item It is not designed to protect against:
    \begin{itemize}
      \item Runtime compromise of the system via a vulnerability
    \end{itemize}
    \item It can be designed to protect against:
    \begin{itemize}
      \item Leaking of some of the cryptographic material
    \end{itemize}
    \item Secure boot is a chain, so its security is a consequence of:
    \begin{itemize}
      \item The security of the root/anchor.
      \item The security of each link, guaranteed by the signature scheme.
    \end{itemize}
  \end{itemize}
\end{frame}

\begin{frame}{Threat Model: the root}
  \begin{itemize}
    \item The root is made of the elements that are initially trusted at boot:
    \begin{itemize}
      \item The hash(es) of the secure boot public key(s)
      \item The bootROM
      \item The CPU
      \item Optionally, harware implementations of cryptography
    \end{itemize}
    \item Of these, only the hash(es) are in an integrator's control
    \item The SoC vendor is very much part of the trusted perimeter.
  \end{itemize}
\end{frame}

\subsection{Example: UEFI secure boot}

\begin{frame}
  \begin{itemize}
    \item Secure boot is part of the UEFI spec
    \item Theoretically support x86-64, ARM32, aarch64, loongarch and RISC-V
    \item Mostly used on x86
    \item UEFI uses a mofified version of Microsoft's Portable Executable (PE) file format
    \item Must be authenticated by the CPU first
    \begin{itemize}
      \item Intel Boot Guard
      \item \href{https://www.ioactive.com/exploring-amd-platform-secure-boot/}{AMD Platform Secure Boot} (PSB)
    \end{itemize}
  \end{itemize}
\end{frame}

\subsection{Example: Raspberry Pi secure boot}

\begin{frame}{Raspberry Pi Secure boot support}
  \begin{itemize}
    \item Since Raspberry Pi4, RPis support secure boot
    \item Raspberry Pi Ltd has 2 whitepapers on the topic:
    \begin{itemize}
      \item \href{https://pip.raspberrypi.com/categories/1260-security/documents/RP-004651-WP/Raspberry-Pi-4-Boot-Security.pdf}{Boot security}
      \item \href{https://pip-assets.raspberrypi.com/categories/1260-security/documents/RP-003466-WP-3-Boot\%20Security\%20Howto.pdf?disposition=inline}{Secure Boot}
    \end{itemize}
    \item The BootROM holds 4 RSA 2048 public keys
    \item These keys are held by the RPi foundation, the first one is a dev key
    \item Broadcom's BCM2711 SoC embeds a
          \href{https://www.raspberrypi.com/documentation/computers/raspberry-pi.html\#one-time-programmable-settings}{one-time programmable (OTP) memory}. It holds:
    \begin{itemize}
      \item revocation status of the ROM keys
      \item Configuration (e.g. boot mode)
      \item 1 slot for a SHA256 hash of an OEM key
      \item 256-bit device unique private key (readable to root in cleartext)
    \end{itemize}
  \end{itemize}
\end{frame}

\begin{frame}{RPi: Secure boot process}
  \begin{itemize}
    \item The RPi BootROM first:
    \begin{itemize}
      \item checks the OTP revocation bits
      \item Verifies \code{bootsys}' signature from EEPROM against a non-revoked key
      \item starts \code{bootsys} (if verification succeeded)
    \end{itemize}
    \item \code{bootsys} then:
    \begin{itemize}
      \item verifies the OEM key in EEPROM against the OTP hash
      \item verifies \code{boot.conf} (in EEPROM) against \code{boot.sig} with the OEM key
      \item verifies \code{bootmain} (in EEPROM) against the embedded hash (which was therefore signed
            with the RPi ROM key)
      \item loads and executes \code{bootmain} if verification succeeded
    \end{itemize}
    \item \code{bootmain} then
    \begin{itemize}
      \item verifies \code{boot.img} against \code{boot.sig} (in boot medium) against the OEM key
      \item loads the \code{boot.img} and starts \code{start.elf}
    \end{itemize}
    \item Finally, \code{start.elf} loads the kernel from the \code{boot.img} ramdisk
  \end{itemize}
\end{frame}

\begin{frame}[fragile]
  \frametitle{RPi: Secure boot overview}
  \begin{itemize}
    \item It is possible to test the \code{boot.img} signature
  \end{itemize}
    \begin{minted}[fontsize=\small]{yaml}
3.04 OTP boardrev b04170 bootrom a a
3.06 Customer key hash 8251a63a2edee9d8f710d63e9da5d639064929ce15a2238986a189ac6fcd3cee
3.13 VC-JTAG unlocked
3.36 RP1_BOOT chip ID: 0x20001927
3.41 bootconf.sig
3.41 hash: f71ede8fad8bea2f853bcff41173ffedde48c5b76ed46bc38fa057ce46e5d58b
3.47 rsa2048: 3f215305d5aff620219da94f6f1294787e3a407102a507da96c28e9195d3ccb2f144cac66919f9d86ba9f54a8d20ff57c80d6d269e6e49a16dc23553974489947fe05bf3b7df5cd2c5040a9eebadca754ff4be50600b06fd9f565639adc859d88052e15e0ff6eecf7fec0386d41f81e5d009b04520bb83f17663b62b1271b9d27ec2344c73a20d42dfd68facd741d48c0453e8149448537abfed1d4805872c16182a3e9f25c0b86e002e88949d62c148a561aa8137c257ce0d3e0ae5761aa64c225e9c9782b2bb613de7d90499567c56218bb18a239d4347967b68b3ebd06eaa48215f16316d2a697bb2e67cb3883068f6284e2ca71d25ce0099a1ceb37a85c9
3.94 RSA verify
3.10 rsa-verify pass (0x0)
    \end{minted}
  \begin{itemize}
    \item On RPi4 , it is also possible to test the EEPROM signature
    \item This is no longer possible on RPi5:
    \begin{itemize}
      \item The EEPROM image must be counter-signed using the OEM key
      \item If the OTP has not been flashed, boot will not proceed
    \end{itemize}
  \end{itemize}
\end{frame}

%\subsection{Secure boot on the Rockchip RK3399}

%\begin{frame}{}
%\end{frame}
\subsection{Example: AHAB on NXP i.MX93}

\begin{frame}{Detailed example: i.MX93}
  \begin{itemize}
    \item i.MX SoCs implement a version of secure boot called High Assurance Boot (HAB)
    \item Starting from i.MX8, moved to
          \href{https://github.com/nxp-imx/uboot-imx/blob/lf_v2025.04/doc/imx/ahab/introduction_ahab.txt}{Advanced HAB} (AHAB)
    \item Uses a table of 4 asymmetric keys called Secure Root Keys (SRKs)
    \item The hash of the SRK table is stored in the SoC's fuses
    \item AHAB actions can be performed using NXP's
          \href{https://spsdk.readthedocs.io/en/latest/index.html\#}{Secure Provisioning SDK (SPSDK)}
  \end{itemize}
\end{frame}

\begin{frame}{The EdgeLock (Secure) Enclave (ELE)}
  \begin{itemize}
    \item The ELE is the security subsystem on the i.MX9 family
    \item Replaces the SEcurity COntroller (SECO) of the i.MX8 family
    \item Also called "Sentinel"
    \item Based on a dedicated RISC-V core
    \item Authenticates all firmware loaded at boot
    \item ELE firmware implements AHAB. Signed using ELE SRKs
          provided as a binary blob by NXP.
  \end{itemize}
\end{frame}

\begin{frame}{The EdgeLock (Secure) Enclave (ELE)}
  \begin{itemize}
    \item Described in the ELE API Reference Guide (IMX93ELEAPI) (NXP account required for download)
    \item HSM capability, API described in \href{https://www.nxp.com/docs/en/reference-manual/RM00284.pdf}{RM00284}
    \begin{itemize}
      \item Key generation
      \item Key storage (no internal NVM, this is important)
      \item Encryption/Decryption
      \item Signature (and verification)
    \end{itemize}
    \item Using the ELE to wrap keys for filesystem encryption is possible
    \item No ELE support for \href{https://docs.kernel.org/security/keys/trusted-encrypted.html}{Trusted Keys}
    \begin{itemize}
      \item LUKS will not be supported, only "naked" cryptsetup
      \item An initramfs will most likely be necessary
      \item Will require some provisioning to initially wrap the key
    \end{itemize}
    \item API accessible from userland via a messaging unit
    \item NXP has a library for this: \href{https://github.com/nxp-imx/imx-secure-enclave}{imx-secure-enclave}
  \end{itemize}
\end{frame}

\begin{frame}{Detailed example: i.MX93 - AHAB container}
  \begin{itemize}
    \item AHAB containers are custom structured binary files
    \item They can be built using SPSDK's nxpimage
    \begin{itemize}
      \item Generate a template: \code{nxpimage ahab get-template}
    \end{itemize}
    \item They will include one or several binaries
    \item The \code{signer} property will indicate how to sign the binaries
  \end{itemize}
\end{frame}

\begin{frame}[fragile]
  \frametitle{Detailed example: i.MX93 - AHAB container template}
  \begin{minted}[fontsize=\scriptsize]{yaml}
# Description: NXP chip family identifier.
family: mimx9352
# --------------===== MCU revision [Optional] =====--------------------------------
# Description: Revision of silicon. The 'latest' name, means most current revision.
# Possible options: <a0, a1, latest>
revision: a1
# --------------===== Memory type [Required] =====----------------------------------
# Description: Specify type of memory used by bootable image description.
memory_type: serial_downloader
init_offset: 0
# --------------===== Primary Image Container Set [Optional] =====-------------------
# Container Set that is validated by ROM and usually contains DDR init and SPL.
# It could be used as pre-prepared binary form of AHAB and also YAML configuration
# file for AHAB. In case that YAML configuration file is used, the Bootable image tool
# builds the AHAB itself.
primary_image_container_set: ahab_primary_container.yaml
secondary_image_container_set: ahab_secondary_container.yaml
  \end{minted}
\end{frame}

\begin{frame}[fragile]
  \frametitle{Detailed example: i.MX93 - AHAB primary container}
  \begin{minted}[fontsize=\scriptsize]{yaml}
family: mimx9352
revision: a1
target_memory: serial_downloader
output: ahab-primary-container-set.bin
containers:
  - binary_container:
      path: mx93a1-ahab-container.img   # ELE firmware, provided and signed by NXP
  - container:
      srk_set: oem
      used_srk_id: 0
      signer: Super_Root_Key_1.pem
      images:
        -
          lpddr_imem_1d: lpddr4_imem_1d_v202201.bin     # LPDDR memory FW in 1D mode
          lpddr_imem_2d: lpddr4_imem_2d_v202201.bin     # LPDDR memory FW in 2D mode
          lpddr_dmem_1d: lpddr4_dmem_1d_v202201.bin     # LPDDR memory data in 1D mode
          lpddr_dmem_2d: lpddr4_dmem_2d_v202201.bin     # LPDDR memory data in 2D mode
          spl_ddr: u-boot-spl.bin       # SPL
      srk_table:
        flag_ca: false
        hash_algorithm: default
        srk_array:
          - Super_Root_Key_1.pub    # 4 lines, one for each SRK
  \end{minted}
\end{frame}

\begin{frame}[fragile]
  \frametitle{Detailed example: i.MX93 - AHAB secondary container}
  \begin{minted}[fontsize=\scriptsize]{yaml}
[...]
containers:
  -
    container:
      srk_set: oem
      used_srk_id: 0
      signer: Super_Root_Key_1.pem
      images:
        - atf: bl31-imx93.bin-optee
        - uboot: u-boot.bin
        - tee: tee.bin
      srk_table:
        flag_ca: false
        hash_algorithm: default
        srk_array:
          - Super_Root_Key_1.pub    # 4 lines, one for each SRK
  \end{minted}
\end{frame}

\begin{frame}
  \frametitle{Detailed example: i.MX93 - AHAB}
  \begin{itemize}
    \item As mentioned, the hashes of the SRKs must be flashed to the SoC's fuses.
    \item The \href{https://github.com/nxp-mcuxpresso/spsdk/blob/master/spsdk/data/devices/mimx9352/fuses.json\#L6638}{SPSDK sources}
          show that those fuses are index \code{0x80} to \code{0x87}
    \item On i.MX93, the hash function being used is SHA256, so the hash will be
          spread over 8 32-bit fuses.
    \item This hash is calculated over all 4 SRKs, they are not fully independent
    \item The fuses will be flashed by the ELE, via the usual message interface
    \begin{itemize}
      \item U-Boot has an \href{https://github.com/nxp-imx/uboot-imx/blob/lf_v2025.04/arch/arm/mach-imx/ele_ahab.c\#L874}{ele\_message} command
    \end{itemize}
    \item This is a sensitive step, so SPSDK can generate a script to automate it
  \end{itemize}
\end{frame}

\begin{frame}[fragile]
  \frametitle{\href{https://github.com/nxp-mcuxpresso/spsdk/blob/master/tests/nxpimage/data/ahab/fuses_scripts/mimx9352_ahab_oem0_srk0_hash_nxpele.bcf}{Example}}
  \begin{minted}[fontsize=\tiny]{yaml}
# nxpele AHAB SRK fuses programming script
# Generated by SPSDK 3.0.0.dev68+g99003d2be
# Family: mimx9352, Revision: latest

# Value: 0xCB2CC774B2DCEC92C840ECA0646B78F8D3661D3A43ED265A490A13ACA75E190A
# Description: SHA256 hash digest of hash of four SRK keys
# Grouped register name: SRKH

# OTP ID: OEM_SRKH7, Value: 0x74C72CCB
write-fuse --index 128 --data 0x74C72CCB
# OTP ID: OEM_SRKH6, Value: 0x92ECDCB2
write-fuse --index 129 --data 0x92ECDCB2
# OTP ID: OEM_SRKH5, Value: 0xA0EC40C8
write-fuse --index 130 --data 0xA0EC40C8
# OTP ID: OEM_SRKH4, Value: 0xF8786B64
write-fuse --index 131 --data 0xF8786B64
# OTP ID: OEM_SRKH3, Value: 0x3A1D66D3
write-fuse --index 132 --data 0x3A1D66D3
# OTP ID: OEM_SRKH2, Value: 0x5A26ED43
write-fuse --index 133 --data 0x5A26ED43
# OTP ID: OEM_SRKH1, Value: 0xAC130A49
write-fuse --index 134 --data 0xAC130A49
# OTP ID: OEM_SRKH0, Value: 0x0A195EA7
write-fuse --index 135 --data 0xA195EA7
  \end{minted}
\end{frame}

\begin{frame}[fragile]
  \frametitle{AHAB status}
  \begin{itemize}
    \item The signature of an AHAB container is always verified
    \item Before the fuses are flashed, this will result in the following error: 
    \begin{minted}[fontsize=\tiny]{yaml}
   > ahab_status 
   Lifecycle: 0x00000008, OEM Open

        0x0287fad6
        IPC = MU APD (0x2)
        CMD = ELE_OEM_CNTN_AUTH_REQ (0x87)
        IND = ELE_BAD_KEY_HASH_FAILURE_IND (0xFA)
        STA = ELE_SUCCESS_IND (0xD6)

        0x0287fad6
        IPC = MU APD (0x2)
        CMD = ELE_OEM_CNTN_AUTH_REQ (0x87)
        IND = ELE_BAD_KEY_HASH_FAILURE_IND (0xFA)
        STA = ELE_SUCCESS_IND (0xD6)
    \end{minted}
    \item Once they are, AHAB should report no events:
    \begin{minted}[fontsize=\tiny]{yaml}
    > ahab_status 
    Lifecycle: 0x00000008, OEM Open


        No Events Found!
    \end{minted}
    \item Then the board can be set to \code{OEM Closed}, and signatures will be enforced
  \end{itemize}
\end{frame}

